% !TeX root = ../../book.tex
\section{皮亚诺公理}
\secbegin{secPeanosAxioms}

本节的目的是忘记我们认为我们所知道的关于自然数的一切，并使用以下基本性质重建我们以前的知识（以及更多！）：

\begin{center}
\textit{每个自然数都可以通过\\从零开始并有限次加一\\这种独特方式获得。}
\end{center}

这有点不太精确 --- 例如，`有限次加一'的含义就不清除。更糟糕的是，我们将在 \Cref{secFiniteSets} 中用自然数来定义`有限'的含义，所以我们最好不要在自然数的定义中引用有限性！

下面的定义准确地捕捉了我们所需的属性，以便描述我们心中 $\mathbb{N}$ 的概念。首先，$\mathbb{N}$ 应该是一个集合。无论这个集合 $\mathbb{N}$ 中的元素实际上\textit{是什么}，我们都会将其视为自然数。其中一个元素尤其需要扮演自然数 $0$ 的角色 --- 这就是\textit{零元素} $z \in \mathbb{N}$；并且需要有一个`加一'的概念 --- 这就是\textit{后继函数} $s : \mathbb{N} \to \mathbb{N}$。因此，给定一个元素 $n \in \mathbb{N}$，它是一个自然数，我们将元素 $s(n)$ 视为自然数`$n+1$'。请注意，这完全是出于直觉的目的：我们将根据 $z$ 和 $s$ 定义 `$+$' 和 `$1$'，反之亦然。

\begin{definition}
\label{defNotionOfNaturalNumbers}
\index{natural numbers!notion of}
\index{Peano's axioms}
\textbf{自然数概念} 是一个集合 $\mathbb{N}$，以及称为\textbf{零元素}的元素 $z \in \mathbb{N}$ 和称为 \textbf{后继函数}的函数 $s : \mathbb{N} \to \mathbb{N}$ ，满足以下性质：
\begin{enumerate}[(i)]
\item $z \not\in s[\mathbb{N}]$；也就是说对于任意 $n \in \mathbb{N}$, $z \ne s(n)$。
\item $s$ 为单射；也就是说，对于所有 $m,n \in \mathbb{N}$，如果 $s(m) = s(n)$，则 $m=n$。
\item $\mathbb{N}$ 由 $z$ 和 $s$ 生成；也就是说，对于所有集合 $X$，如果 $z \in X$ 且对于所有 $n \in \mathbb{N}$ 我们有 $n \in X \Rightarrow s(n) \in X$，则 $\mathbb{N} \subseteq X$.
\end{enumerate}
性质 (i), (ii) 和 (iii) 称为 \textbf{皮亚诺公理}。
\end{definition}

请注意，\Cref{defNotionOfNaturalNumbers} 并未指明 $\mathbb{N}$、$z$ 和 $s$ 实际上是什么；它只是指明它们必须满足的性质。事实证明，我们使用什么自然数概念并不重要，因为任何两个概念本质上都是相同的。我们不必担心这里的具体细节 --- 该任务留给 \Cref{secConstructions}：自然数的特定概念在 \Cref{cnsNaturalNumbersVonNeumann} 中定义，并且所有自然数概念`本质上都是相同的'在 \Cref{thmNNNUnique} 中给出精确证明。

只需从元素 $z \in \mathbb{N}$ 和后继函数 $s : \mathbb{N} \to \mathbb{N}$ 出发，我们可以定义所有我们熟悉的有关自然数的概念，并证明所有我们认为理所当然的性质。

比如，我们定义 $z$ 表示 `$0$'，定义 $s(z)$ 表示 `$1$'，定义 $s(s(z))$ 表示 `$2$'，以此类推。所以 `$12$' 可以定义为
\[ s(s(s(s(s(s(s(s(s(s(s(s(z)))))))))))) \]

从现在开始，$\mathbb{N}$ 的零元素就写成 $0$ 而不是 $z$。如果我们能将 $s(n)$ 写成 `$n+1$' 那就太好了，但我们必须首先定义 `$+$' 的含义。为了做到这一点，我们需要一种涉及自然数的定义表达方式；这就是 \textit{递归定理} 允许我们做的事情。

\begin{theorem}[递归定理]
\label{thmRecursion}
\index{recursion!recursion theorem}
设 $X$ 为集合。对于所有 $a \in X$ 和所有 $h : \mathbb{N} \times X \to X$，存在唯一一个函数 $f : \mathbb{N} \to X$ 使得对于所有 $n \in \mathbb{N}$, $f(0) = a$ 且 $f(s(n)) = h(n, f(n))$。
\end{theorem}

\begin{cproof}
设 $a \in X$ 且 $h : \mathbb{N} \times X \to X$。我们需要分别证明 $f$ 的存在性和唯一性。

\begin{itemize}
\item 通过指定 $f(0) = a$ 和 $f(s(n)) = h(n, f(n))$ 定义 $f : \mathbb{N} \to X$。因为 $h$ 为函数且 $s$ 为单射，只要给出 $f(n)$ 的定义，就可以保证 $x \in X$ 的存在性和唯一性，使得$f(n) = x$，因此我们只需要验证完整性即可。

所以设 $D = \{ n \in \mathbb{N} \mid f(n) \text{ is defined} \}$。则：

\begin{itemize}
\item $0 \in D$，因为 $f(0)$ 被定义为等于 $a$。
\item 设 $n \in \mathbb{N}$ 并假设 $n \in D$。则 $f(n)$ 被定义且 $f(s(n)) = h(n, f(n))$，所以 $f(s(n))$ 被定义，因此 $s(n) \in D$。
\end{itemize}

%% BEGIN EXTRACT (xtrConclusionExample) %%
根据 \Cref{defNotionOfNaturalNumbers} 的性质 (iii)，我们有 $\mathbb{N} \subseteq D$，因此 $f(n)$ 对所有 $n \in \mathbb{N}$ 有定义。
%% END EXTRACT %%

\item 要证明 $f$ 是唯一的，假设 $g : \mathbb{N} \to X$ 为另一个函数对所有 $n \in \mathbb{N}$ 满足 $g(0) = a$ 且 $g(s(n)) = h(n, g(n))$。

要想得到 $f = g$，设 $E = \{ n \in \mathbb{N} \mid f(n) = g(n) \}$。则
\begin{itemize}
\item $0 \in E$，因为 $f(0) = a = g(0)$。
\item 设 $n \in \mathbb{N}$ 并假设 $n \in E$。则 $f(n) = g(n)$，因此
\[ f(s(n)) = h(n,f(n)) = h(n,g(n)) = g(s(n)) \]
所以 $s(n) \in E$。
\end{itemize}
再次满足 \Cref{defNotionOfNaturalNumbers} 的性质(iii)，因此 $\mathbb{N}\subseteq E$。由此可见，对于所有 $n \in \mathbb{N}$, $f(n) = g(n)$，因此 $f=g$。
\end{itemize}

因此我们确立了对于所有 $n \in \mathbb{N}$ 满足 $f(0) = a$ 且 $f(s(n)) = h(n, f(n))$ 的函数 $f : \mathbb{N} \to X$ 的存在性和唯一性。
\end{cproof}

递归定理允许我们通过\textit{递归}来定义涉及自然数的表达式；这就是 \Cref{strDefinitionByRecursion}。

\begin{strategy}[通过递归定义]
\label{strDefinitionByRecursion}
\index{recursion!definition by recursion}
要想定义函数 $f : \mathbb{N} \to X$，只需定义 $f(0)$，并且对于给定的 $n \in \mathbb{N}$，假设 $f(n)$ 已定义，并根据 $n$ 和 $f(n)$ 定义 $f(s(n))$ 即可。
\end{strategy}

\begin{example}
我们可以使用递归来定义自然数的加法，如下所示。

对于给定的 $m \in \mathbb{N}$，我们可以通过如下递归定义函数 $\mathrm{add}_m : \mathbb{N} \to \mathbb{N}$：
\[ \mathrm{add}_m(0) = m \quad \text{且} \quad \mathrm{add}_m(s(n)) = s(\mathrm{add}_m(n)) \text{ 对于所有 } n \in \mathbb{N} \]
用更熟悉的方法表示，对于 $m,n \in \mathbb{N}$，定义表达式 `$m+n$' 表示 $\mathrm{add}_m(n)$。表达 $\mathrm{add}_m(n)$ 递归定义的另一种方式是，对于每个 $m \in \mathbb{N}$，我们通过在 $n$ 上递归定义 $m+n$ 如下：
\[ m+0 = m \quad \text{且} \quad m+s(n) = s(m+n) \text{ 对于所有 } n \in \mathbb{N} \]
\end{example}

我们可以使用加法的递归定义来证明数字之间熟悉的等式关系。下面是命题 $2+2=4$ 的证明。这可能看起来很愚蠢，但请注意表达式`$2+2=4$'实际上是以下内容的简写：
\[ \mathrm{add}_{s(s(0))} (s(s(0))) = s(s(s(s(0)))) \]
因此，我们必须小心地在其证明中应用定义。

\begin{proposition}
\label{propTwoPlusTwoEqualsFour}
$2+2=4$
\end{proposition}

\begin{cproof}
我们使用加法的递归定义。
\begin{align*}
2 + 2 &= 2 + s(1) && \text{因为 $2=s(1)$} \\
&= s(2+1) && \text{根据 $+$ 的定义} \\
&= s(2+s(0)) && \text{因为 $1=s(0)$} \\
&= s(s(2+0)) && \text{根据 $+$ 的定义} \\
&= s(s(2)) && \text{根据 $+$ 的定义} \\
&= s(3) && \text{应为 $3=s(2)$} \\
&= 4 && \text{应为 $4=s(3)$}
\end{align*}
\end{cproof}

下面的结果允许我们删除符号 `$s(n)$' 而只写 `$n+1$'。

\begin{proposition}
\label{propSuccessorIsPlusOne}
对于所有 $n \in \mathbb{N}$，我们有 $s(n) = n+1$。
\end{proposition}

\begin{cproof}
设 $n \in \mathbb{N}$。根据加法的递归定义，我们有
\[ n+1 = n+s(0) = s(n+0) = s(n) \]
\end{cproof}

根据 \Cref{propSuccessorIsPlusOne}，我们现在将放弃使用 $s(n)$ 符号，并改用 $n+1$。

我们也可以通过递归来定义乘法和幂运算。

\begin{example}
给定 $m \in \mathbb{N}$。对于所有 $n \in \mathbb{N}$ 递归定义 $m \cdot n$ 如下：
\[ m \cdot 0 = 0 \quad \text{且} \quad m \cdot (n+1) = (m \cdot n) + m \text{ 对于所有 } n \in \mathbb{N} \]
形式上，我们所做的是通过 $\mathrm{mult}_m(z)=z$ 和 $\mathrm{mult}_m(s(n)) = \mathrm{add}_{\mathrm{mult}_m(n)}(m)$ 对于所有 $n \in \mathbb{N}$ 递归地定义一个函数  $\mathrm{mult}_m : \mathbb{N} \to \mathbb{N}$。但我们提供的定义更容易理解。
\end{example}

\begin{proposition}
\label{propTwoTimesTwoEqualsFour}
$2 \cdot 2 = 4$
\end{proposition}

\begin{cproof}
我们使用加法和递归的递归定义。
\begin{align*}
2 \cdot 2 &= 2 \cdot (1+1) && \text{因为 $2=1+1$} \\
&= (2 \cdot 1) + 2 && \text{根据 $\cdot$ 的定义} \\
&= (2 \cdot (0+1)) + 2 && \text{因为 $1 = 0 + 1$} \\
&= ((2 \cdot 0) + 2) + 2 && \text{根据 $\cdot 的定义$} \\
&= (0+2) + 2 && \text{根据 $\cdot$ 的定义} \\
&= (0+(1+1)) + 2 && \text{因为 $2=1+1$} \\
&= ((0+1)+1) + 2 && \text{根据 $+$ 的定义} \\
&= (1+1) + 2 && \text{因为 $1=0+1$} \\
&= 2+2 && \text{因为 $2=1+1$} \\
&= 4 && \text{根据 \Cref{propTwoPlusTwoEqualsFour}}
\end{align*}
\end{cproof}

\begin{exercise}
给定 $m \in \mathbb{N}$，对于所有 $n \in \mathbb{N}$ 递归定义 $m^n$，并证明使用幂、乘法和加法的递归定义证明 $2^2 = 4$。
\end{exercise}

我们可以用尽余生来进行涉及递归定义的算术运算，因此我们将就此打住，回到我们所知的理所当然正确的算术运算法则。

然而，我们还需要通过递归来定义更多概念，以便稍后在证明中使用。

\begin{definition}
\label{defSumOfRealNumbers}
给定 $n \in \mathbb{N}$, $n$ 个实数 $a_1, a_2, \dots, a_n$ 之\textbf{和} 为实数 $\sum_{k=1}^n a_k$ 在 $n \in \mathbb{N}$ 上的递归定义为
\[ \sum_{k=1}^0 a_k = 0 \quad \text{且} \quad \sum_{k=1}^{n+1} a_k = \left( \sum_{k=1}^n a_k \right) + a_{n+1} \text{ 对于所有 } n \in \mathbb{N} \]
\end{definition}

\begin{definition}
\label{defProductOfRealNumbers}
给定 $n \in \mathbb{N}$, $n$ 个实数 $a_1, a_2, \dots, a_n$ 之\textbf{积} 为实数 $\prod_{k=1}^n a_k$ 在 $n \in \mathbb{N}$ 上的递归定义为
\[ \prod_{k=1}^0 a_k = 1 \quad \text{且} \quad \prod_{k=1}^{n+1} a_k = \left( \prod_{k=1}^n a_k \right) \cdot a_{n+1} \text{ 对于所有 } n \in \mathbb{N} \]
\end{definition}

\begin{example}
设对于所有 $i \in \mathbb{N}$, $x_i=i^2$。则
\[ \sum_{i=1}^5 x_i = 1 + 4 + 9 + 16 + 25 = 55 \]
且
\[ \prod_{i=1}^5 x_i = 1 \cdot 4 \cdot 9 \cdot 16 \cdot 25 = 14400 \]
\end{example}

\begin{exercise}
设 $x_1, x_2 \in \mathbb{R}$。严格根据索引和和索引积的定义，证明
\[ \sum_{i=1}^2 x_i = x_1 + x_2 \quad \text{且} \quad \prod_{i=1}^2 x_i = x_1 \cdot x_2 \]
\end{exercise}

\subsection*{二项式和阶乘}

\begin{definition}[在 \Cref{defFactorial} 中重新定义]
\label{defFactorialRecursive}
\index{factorial}
设 $n \in \mathbb{N}$。$n$ 的\textbf{阶乘}，写做 $n!$\nindex{nfactorial}{$n"!$}{factorial}，递归定义为
\[ 0! = 1 \quad \text{且} \quad (n+1)! = (n+1) \cdot n! \text{ 对于所有 } n \ge 0 \]
\end{definition}

换句话说，我们有
\[ n! = \prod_{i=1}^n i \]
对于所有 $n \in \mathbb{N}$ --- 回想一下 \Cref{defProductOfRealNumbers} 看看为什么这些定义实际上只是同一事物的两种表述。

\begin{definition}[在 \Cref{defBinomialCoefficient} 中重新定义]
\label{defBinomialCoefficientRecursive}
\index{binomial coefficient}
\nindex{nchoosek}{$\binom{n}{k}$}{binomial coefficient}
设 $n,k \in \mathbb{N}$。\textbf{二项式系数} $\binom{n}{k}$ \inlatex{binom\{n\}\{k\}}\lindexmmc{binom}{$\binom{n}{k}$} (读作 `$n$ 选 $k$') 在 $n$ \textit{和} $k$ 上递归定义为
\[ \binom{n}{0}=1, \quad \binom{0}{k+1} = 0, \quad \binom{n+1}{k+1} = \binom{n}{k} + \binom{n}{k+1} \]
\end{definition}

这个定义产生了一种计算二项式系数的算法：它与所谓的\textbf{帕斯卡三角}\index{Pascal's triangle}相匹配，每个二项式系数为位于其上方的两个系数之和（忽略零）：

\begin{center}\begin{tabular}{ccc}
$\binom{0}{0}$ && $1$ \\
$\binom{1}{0}$ \quad $\binom{1}{1}$ && $1$ \quad $1$ \\
$\binom{2}{0}$ \quad $\binom{2}{1}$ \quad $\binom{2}{2}$ & = & $1$ \quad $2$ \quad $1$ \\
$\binom{3}{0}$ \quad $\binom{3}{1}$ \quad $\binom{3}{2}$ \quad $\binom{3}{3}$ && $1$ \quad $3$ \quad $3$ \quad $1$ \\
$\binom{4}{0}$ \quad $\binom{4}{1}$ \quad $\binom{4}{2}$ \quad $\binom{4}{3}$ \quad $\binom{4}{4}$ && $1$ \quad $4$ \quad $6$ \quad $4$ \quad $1$ \\
$\binom{5}{0}$ \quad $\binom{5}{1}$ \quad $\binom{5}{2}$ \quad $\binom{5}{3}$ \quad $\binom{5}{4}$ \quad $\binom{5}{5}$ && $1$ \quad $5$ \quad $10$ \quad $10$ \quad $5$ \quad $1$ \\
$\vdots$ \qquad $\vdots$ \qquad $\vdots$ \qquad $\vdots$ \qquad $\vdots$ && $\vdots$ \qquad $\vdots$ \qquad $\vdots$ \qquad $\vdots$
\end{tabular}\end{center}

\begin{exercise}
写出帕斯卡三角接下来的两行。
\end{exercise}